% Chapter 6: Ridders' extrapolation (methods author). PLAN.md section 1.4; sidebars S10, S11 and S12; figures F12
% (table*, with its companion diagram), F13, F14 (figure*) with F14b (margin), F15 (margin); examples E9 and E10.
% Style decision 19: the margin floats F14b and F15 are kept a page away from the full-width F12 and F14.
% Revised at stage 6 per reviews/R1-decision.md (T7, T12, T13, T15, numerics-10, 12, 14, 15, 16, 17c, 17i, 17j,
% 17k, library-14, engineer-12, 19, 22, 24, evidence-09, 10, 17); the changes are listed in responses/R1-06-ridders.md.
\chapter{Ridders' extrapolation}\label{ch:ridders}

Every stencil of the last three chapters returns a number and nothing else. Its step is a guess, the default or
the user's, and whether the guess was good shows only when the answer is known from elsewhere. Ridders'
method,\cite{ridders1982} in the form the library implements it, removes the guess and adds the one thing a
stencil cannot give: an estimate of its own error. It evaluates the central difference at a step, halves the step,
and again, extrapolating the sequence to the limit of zero step until rounding stops the improvement; the size of
the last improvement is the estimate. The library offers it as one more method tag,
\code{derivative<Ridders>(fn, x)}, for every derivative order. This chapter explains the extrapolation, shows the
tableau on the nitrogen example, says how to read \code{abs_error} and when not to trust it, what it all costs, and
what its one-sided forms do at the edge of a domain.

\section{Richardson's idea}\label{sec:ridders-richardson}

The central difference $\Dc{h} = (f(x + h) - f(x - h)) / 2h$ has an error expansion in even powers of the step
(\cref{ch:taylor}, \cref{sb:remainder}):
\[
  \Dc{h} = f'(x) + a_2 h^2 + a_4 h^4 + a_6 h^6 + \cdots ,
\]
with coefficients $a_{2m}$ that depend on $f$ but not on $h$. Halving the step divides the leading error term by
4, so $\Dc{h}$ and $\Dc{h/2}$ are two measurements with a known relation between their errors, and one
combination cancels the term:
\begin{equation}
  \frac{4 \Dc{h/2} - \Dc{h}}{3} = f'(x) - \frac{a_4}{4} h^4 + \cdots .
  \label{eq:ridders-richardson}
\end{equation}
This is Richardson extrapolation,\cite{richardson1911,richardson1927} the deferred approach to the limit: from two
estimates of order $h^2$, one of order $h^4$, without knowing $a_2$. Applied once at a fixed step it is nothing
new; \cref{sb:richardson-disguise} showed that \cref{eq:ridders-richardson} is the four-point rule
\code{Central<4>} written differently. What makes it a method is repetition: the new estimate has the same form,
an expansion in even powers with the leading term now $h^4$, so a second combination with the factor
$4^2 = 16$ removes that term, a third with $4^3$ the next, and so on, each from values the previous stage already
has. Burden et al.\cite{burden2016} (Section 4.2) present the tableau this builds and Sidi\cite{sidi2003} the general
theory. Ridders' paper supplies the idea; the rules, the ratio 2, a rounding floor, a two-level stall, a final
check and restarts, are the library's own, chosen on a prototype (\cref{sb:floor-coverage}), and every statistic
in this chapter is of them.

\section{The tableau in one picture}\label{sec:ridders-tableau}

\begin{figure}
\sidecaption{The shape of the tableau: column 0 holds the base differences $D_i$ at the halved steps; every other
entry is built from its left neighbour and that neighbour's parent one level up. The diagonal entries are the latest
of each level. A diagram without data.\label{fig:ridders-tableau-diagram}}
\centering
% F12 companion (Chapter 6, text width): the shape of Ridders' tableau as a diagram (no data). Column 0 holds the
% base differences D_i at the halved steps h_0 / 2^i; every other entry is built from its left neighbour and that
% neighbour's parent one level up, in Burden's difference form A_{i,j} = A_{i,j-1} + (A_{i,j-1} - A_{i-1,j-1}) /
% (4^j - 1). The one entry in accent shows the two parents it depends on; the diagonal entries (the latest entry
% of each level, against which the returned entry is checked) are in the darker ink.
\begin{tikzpicture}[x=14mm, y=7mm, font=\footnotesize,
  entry/.style={inner sep=1.5pt, text=black!55},
  diagonal/.style={inner sep=1.5pt, text=black!85},
  live/.style={inner sep=1.5pt, text=accent},
  dep/.style={-{Stealth[length=3pt]}, accent, thin}]
  % the entries: level i downwards, column j to the right; the diagonal darker
  \foreach \i in {0,...,4}
    \foreach \j in {0,...,\i} {
      \ifnum\i=\j
        \node[diagonal] (A\i\j) at (\j, -\i) {$\Aij{\i}{\j}$};
      \else
        \node[entry] (A\i\j) at (\j, -\i) {$\Aij{\i}{\j}$};
      \fi
    }
  % the one to notice and its two parents
  \node[live] (A32) at (2, -3) {$\Aij{3}{2}$};
  \node[live] (A31) at (1, -3) {$\Aij{3}{1}$};
  \node[live] (A21) at (1, -2) {$\Aij{2}{1}$};
  \draw[dep] (A31) -- (A32);
  \draw[dep] (A21) -- (A32);
  \node[accent, anchor=north, font=\scriptsize] at (2, -4.6)
    {$\Aij{3}{2} = \Aij{3}{1} + \dfrac{\Aij{3}{1} - \Aij{2}{1}}{4^2 - 1}$};
  % the steps of the levels
  \foreach \i/\divisor in {0/1, 1/2, 2/4, 3/8, 4/16}
    \node[anchor=east, text=black!70, font=\scriptsize] at (-0.4, -\i)
      {level $\i$, step \ifnum\i=0 $h_0$\else $h_0 / \divisor$\fi};
  % column headings
  \node[anchor=south, text=black!70, font=\scriptsize, align=center] at (0, 0.45) {column 0\\(the base $D_i$)};
  \foreach \j in {1,...,4}
    \node[anchor=south, text=black!70, font=\scriptsize] at (\j, 0.45) {column \j};
  % the diagonal
  \node[anchor=west, text=black!85, font=\scriptsize, align=left] at (4.35, -3.5) {the diagonal $\Aij{i}{i}$:\\the latest entry\\of each level};
\end{tikzpicture}

\end{figure}

Write $h_i = h_0 / 2^i$ for the step of level $i$ and $D_i = \Dc{h_i}$ for the base difference there. The
extrapolated values form a triangular tableau (\cref{fig:ridders-tableau-diagram}), one row per level and one
column per cancelled term:
\begin{equation}
  \Aij{i}{0} = D_i, \qquad
  \Aij{i}{j} = \Aij{i}{j-1} + \frac{\Aij{i}{j-1} - \Aij{i-1}{j-1}}{4^j - 1}, \quad 1 \le j \le i .
  \label{eq:ridders-tableau}
\end{equation}
Column $j$ is free of the terms up to $h^{2j}$, and the entry $\Aij{i}{j}$ is built from its left neighbour and
that neighbour's parent one level up. The form of \cref{eq:ridders-tableau}, a correction added to the left
neighbour rather than the weighted combination of \cref{eq:ridders-richardson}, is the one Burden et al. write
and the library computes, for reasons given in \cref{sb:difference-form}; the two are algebraically the same.

\begin{table*}
\caption{Ridders' tableau for $\cp$ at $T_0$ from $h_0 = T_0 / 32$: the \num{\valTableauLevelsUsed} levels the method
used, the entries $\Aij{i}{j}$ above and their relative errors with the rounding floor $r_i$ below; the entry
Ridders returned (level \num{\valTableauBestLevel}, column \num{\valTableauBestColumn}) in the accent colour.\label{tab:ridders-tableau}}
\centering
% F12 tab:ridders-tableau (Chapter 6). Ridders' tableau for cp at T0 from h_0 = T0 / 32: the levels the method
% used, the entries A_{i,j} (first table) and their relative errors with the rounding floor r_i (second table).
% Both bodies are generated (DS8, data/table-ridders-tableau.tex and data/table-ridders-tableau-errors.tex, seven
% columns: \valTableauLevelsUsed); the entry Ridders returned is flagged in accent by the generator. Needs the full
% width (table*): nine and ten columns of twelve-digit numbers do not fit the 114 mm text block.
% The bodies are read with the editor's \tablebody (macros.tex), the TeX primitive \input: LaTeX's \input{...} runs
% file hooks after the file, whose tokens start a new table row before \bottomrule and give "Misplaced \noalign".
\footnotesize\setlength{\tabcolsep}{3.5pt}
\begin{tabular}{@{}rr*{7}{r}@{}}
\toprule
 & & \multicolumn{7}{c}{entry $\Aij{i}{j}$ in J/(mol K)} \\
\cmidrule(lr){3-9}
level $i$ & $h_i$ (K) & $j = 0$ & $j = 1$ & $j = 2$ & $j = 3$ & $j = 4$ & $j = 5$ & $j = 6$ \\
\midrule
\tablebody{ridders-tableau}
\bottomrule
\end{tabular}
\par\medskip
\begin{tabular}{@{}rr*{7}{r}r@{}}
\toprule
 & & \multicolumn{7}{c}{relative error $\abs{\Aij{i}{j} - \cp} / \cp$} & floor \\
\cmidrule(lr){3-9}
level $i$ & $h_i$ (K) & $j = 0$ & $j = 1$ & $j = 2$ & $j = 3$ & $j = 4$ & $j = 5$ & $j = 6$ & $r_i / \cp$ \\
\midrule
\tablebody{ridders-tableau-errors}
\bottomrule
\end{tabular}

\end{table*}

\Cref{tab:ridders-tableau} is the tableau for the nitrogen example, $\cp = \dd \enth / \dd T$ at $T_0$, from the default
initial step $h_0 = T_0 / 32$, about \SI{\valCpRiddersInitialStep}{\kelvin}. Read it by columns. Column 0 is the
plain central difference, which at a step of thirty kelvin is \num{\valTableauBaseFirstRelError} off and gains a
factor 4 per level, as \cref{ch:taylor} said it would. Column 1 is already \code{Central<4>} in disguise and is
far more accurate from its first entry on; column 2 is at the rounding floor by level 3, and from there on the
entries of columns 2 and beyond agree with each other and with the closed form to the last digit shown. The lower
table shows what the upper cannot: the relative errors, and in the last column the rounding floor $r_i$ of each
level, the error that the rounding of the two samples alone imposes on $D_i$, which doubles per level as the step
halves. The errors of the extrapolated entries stop falling once they reach the floor, and from then on they drift
up with it.

\begin{figure}
\sidecaption{Relative error of $\Aij{i}{j}$ against the level $i$ over 16 levels, one curve per column $j$, the
rounding floor $r_i$ dashed: column 0 descends slowly while each higher column reaches the floor within a few
levels; the floor doubles per level, so the entries then drift up. Ridders' stop and its returned entry are
marked. The extrapolation reaches in a few levels what the base difference would need a dozen for; after that,
more levels only add rounding.\label{fig:ridders-convergence}}
\centering
% F13 fig:ridders-convergence (Chapter 6, text width). Relative error of the tableau entries A_{i,j} against the
% level i over 16 levels for cp at T0 (DS8, data/ridders-tableau.csv, long form: one row per (i, j); each curve
% selects its column j with a y expression and discards the other rows). Column 0 (the base difference) descends
% slowly; each higher column reaches the rounding floor r_i (dashed) within a few levels; the floor doubles per
% level, so the entries then drift up. Ridders' stop and its returned entry come from \valTableauStopLevel,
% \valTableauBestLevel and \valTableauBestRelError (the entry matched bit for bit by the generator).
\begin{tikzpicture}
\begin{axis}[eqtufte, width=95mm, height=48mm,
  ymode=log, log basis y=10, xmin=0, xmax=15.5, ymin=1e-16, ymax=1e-4,
  xtick={0,2,...,14}, ytick={1e-16,1e-14,1e-12,1e-10,1e-8,1e-6,1e-4},
  xlabel={level $i$ (step $h_0 / 2^i$, $h_0 = T_0/32$)}, ylabel={relative error of $\Aij{i}{j}$},
  unbounded coords=discard]
  % the rounding floor r_i (one value per level, taken from the j = 0 rows)
  \addplot[muted, dashed] table[col sep=comma, comment chars={\#}, x=i,
    y expr={\thisrow{j} == 0 ? \thisrow{floor} : nan}] {data/ridders-tableau.csv};
  \node[muted, font=\scriptsize, anchor=south, inner sep=1pt] at (axis cs:12.5, 3.5e-11) {floor $r_i$};
  % columns 0 to 4
  \addplot[black!85, mark=*, mark size=1pt] table[col sep=comma, comment chars={\#}, x=i,
    y expr={\thisrow{j} == 0 ? \thisrow{error} : nan}] {data/ridders-tableau.csv};
  \addplot[black!70, mark=square*, mark size=1pt] table[col sep=comma, comment chars={\#}, x=i,
    y expr={\thisrow{j} == 1 ? \thisrow{error} : nan}] {data/ridders-tableau.csv};
  \addplot[black!55, mark=triangle*, mark size=1.2pt] table[col sep=comma, comment chars={\#}, x=i,
    y expr={\thisrow{j} == 2 ? \thisrow{error} : nan}] {data/ridders-tableau.csv};
  \addplot[black!45, mark=diamond*, mark size=1.2pt] table[col sep=comma, comment chars={\#}, x=i,
    y expr={\thisrow{j} == 3 ? \thisrow{error} : nan}] {data/ridders-tableau.csv};
  \addplot[black!35, mark=pentagon*, mark size=1.2pt] table[col sep=comma, comment chars={\#}, x=i,
    y expr={\thisrow{j} == 4 ? \thisrow{error} : nan}] {data/ridders-tableau.csv};
  \node[font=\scriptsize, anchor=south west, inner sep=1pt] at (axis cs:7.3, 6e-9) {$j = 0$: the base difference $D_i$};
  \node[font=\scriptsize, anchor=south west, inner sep=1pt] at (axis cs:3.2, 3e-11) {$j = 1$};
  \node[font=\scriptsize, anchor=north west, inner sep=1pt] at (axis cs:6.3, 7e-16) {$j = 2, 3, 4$ (they coincide)};
  % Ridders' stop and the entry it returned
  \draw[accent, thin, dashed] (axis cs:\valTableauStopLevel, 1e-16) -- (axis cs:\valTableauStopLevel, 1e-4);
  \node[accent, font=\scriptsize, anchor=south west, inner sep=1pt] at (axis cs:\valTableauStopLevel, 3e-5) {Ridders stops};
  \node[accent, draw, circle, thin, inner sep=0pt, minimum size=6pt] at (axis cs:\valTableauBestLevel, \valTableauBestRelError) {};
  \node[accent, font=\scriptsize, anchor=south east, inner sep=1pt] at (axis cs:3.85, 1.1e-13) {returned};
\end{axis}
\end{tikzpicture}

\end{figure}

\Cref{fig:ridders-convergence} continues the same tableau to sixteen levels, beyond where Ridders stopped, so that
the whole picture is visible. The base difference alone needs \num{\valTableauBaseBestLevel} halvings to reach its
best, \num{\valTableauBaseBestRelError}, exactly the V of \cref{ch:rounding} traced level by level; column 2 is at
the floor by level 3. Ridders' method, measured on this call, used \num{\valTableauLevelsUsed} levels,
\num{\valCpRiddersEvals} evaluations in all, and returned the entry $A_{\num{\valTableauBestLevel},\num{\valTableauBestColumn}}$,
whose relative error is \num{\valTableauBestRelError}, for the cost of \num{\valTableauLevelsUsed} central
differences, where \code{Central<2>} with its default step gave \num{\valCpDefaultRelError} for one. The example
that makes this call counts the evaluations with a wrapper around \code{fn}:
\snippet{06-ridders-basic}{ridders-cp}

\section{Reading \texttt{abs\_error}}\label{sec:ridders-estimate}

Where does the number after $\pm$ come from, and what does it mean? Every extrapolated entry of the tableau
carries an estimate of its own error,
\begin{equation}
  e_{i,j} = \max\bigl( \abs{\Aij{i}{j} - \Aij{i-1}{j-1}},\ r_i \bigr), \qquad j \ge 1,
  \label{eq:ridders-estimate}
\end{equation}
the change of the entry from its parent, but never less than the rounding floor $r_i$ of its level. The change
is the textbook estimate of successive approximations: the parent's error is $4^j$ times that of the entry's left
neighbour, and the entry's own error is smaller still, so the difference between entry and parent is nearly the
parent's whole error, which overestimates the entry's (\cref{sec:derivations-richardson}). The floor is what keeps
the estimate honest once truncation is gone: the entries of a converged column differ by rounding alone, the
smallest of several such differences is small by luck rather than by accuracy, and without the floor the method
would pick it and report an error it cannot vouch for. The floor of level $i$ is the rounding error of the base
difference, $(\epsmach \abs{f(x - h_i)} + \epsmach \abs{f(x + h_i)})$ divided by the distance between the two
points, under the assumption that each value of \code{fn} is accurate to about $\epsmach$.

The method keeps the entry with the smallest estimate, and stops when two further levels produce no smaller one,
because each level doubles the floor, so more levels only add rounding (on the nitrogen example, the last two levels of
\cref{tab:ridders-tableau}). The returned \code{abs_error} is that smallest estimate, raised if necessary to the
distance between the returned entry and the latest diagonal entry $\Aij{i}{i}$, which catches an entry that only
looked good. On the nitrogen example the estimate is \SI{\valCpRiddersEstimate}{\joule\per\mole\per\kelvin} against an
actual error of \SI{\valCpRiddersActual}{\joule\per\mole\per\kelvin}; note that the entries of the last level are
more accurate than the one returned, but their estimates, pinned at a larger floor, cannot say so, and the
method reports what it can vouch for rather than what happens to be best. The exact stopping rules, the restart
rules and the constants are in \cref{sb:floor-coverage}, with the evidence that they hold.

\begin{figure*}
\centering
% F14 fig:ridders-coverage (Chapter 6, full width, figure*). Ridders' error estimate against the actual error over
% the random study (DS9): N = 1 in dark grey, N = 2 muted; absolute values (data/README.md: relative ones blow up
% near a zero of f^(N)). The diagonal (accent) is estimate = actual error: a point above it is covered. An actual
% error of exactly 0 (an exact stop) cannot be shown on a log axis and is clipped to eps |f^(N)| with
% eps = 2^-52; the caption counts those cases (\valRiddersZeroErrorFirst, \valRiddersZeroErrorSecond).
% The plot reads data/ridders-coverage-plot-N{1,2}.csv, the two columns it needs with the clipping already applied
% by the generator: with the full ten-column files and an x expression per point, pgfplots used 2.6 of TeX's
% 5 million words of memory for this one figure and the book did not build (editor, stage 2, decision S2-17).
\begin{tikzpicture}
\begin{axis}[eqtufte, width=150mm, height=68mm,
  xmode=log, ymode=log, log basis x=10, log basis y=10,
  xmin=1e-19, xmax=1e-8, ymin=1e-18, ymax=1e-8,
  xtick={1e-18,1e-16,1e-14,1e-12,1e-10,1e-8}, ytick={1e-18,1e-16,1e-14,1e-12,1e-10,1e-8},
  xlabel={actual error $\abs{\text{value} - f^{(\dorder)}(x)}$}, ylabel={error estimate \code{abs_error}}]
  \addplot[only marks, mark=*, mark size=0.5pt, black!35] table[col sep=comma, comment chars={\#},
    x=actual, y=estimate] {data/ridders-coverage-plot-N2.csv};
  \addplot[only marks, mark=*, mark size=0.5pt, black!80] table[col sep=comma, comment chars={\#},
    x=actual, y=estimate] {data/ridders-coverage-plot-N1.csv};
  \addplot[accent, thin, domain=1e-18:1e-8, samples=2] {x};
  \node[accent, font=\scriptsize, anchor=north west, inner sep=2pt, rotate=24] at (axis cs:2e-17, 2e-17) {estimate $=$ actual error};
  \node[font=\scriptsize, anchor=west, inner sep=1pt] at (axis cs:1.5e-18, 1e-13) {$\dorder = 1$};
  \node[font=\scriptsize, anchor=west, inner sep=1pt, text=black!55] at (axis cs:1.5e-18, 1e-10) {$\dorder = 2$ (muted)};
  \node[font=\scriptsize, anchor=south east, inner sep=1pt, text=black!70] at (axis cs:1e-8, 1.3e-18) {covered: above the diagonal};
\end{axis}
\end{tikzpicture}

\caption{Ridders' error estimate against the actual error over the random study ($\dorder = 1$ dark,
$\dorder = 2$ muted; \num{\valRiddersFunctions} smooth functions on $[0.2, 3]$, \num{\valRiddersPointsPerFunction} points
each, no noise; no singularity within $h_0$ for $f'$, while for $f''$ the initial step $\scale / 4$ reaches past
zero below $x = 0.25$, the restarts of \cref{sb:floor-coverage}), the diagonal drawn: points above it are covered,
\SI{\valRiddersCoverageFirst}{\percent} of the \num{\valRiddersCases} cases for $\dorder = 1$ and
\SI{\valRiddersCoverageSecond}{\percent} for $\dorder = 2$. Every third case is drawn (\num{\valRiddersPlotCases} per
order; all enter the statistics). An actual error of exactly 0 (\num{\valRiddersZeroErrorFirst} and
\num{\valRiddersZeroErrorSecond} cases) is clipped to $\epsmach \abs{f^{(\dorder)}}$.\label{fig:ridders-coverage}}
\end{figure*}

How good is the estimate? \Cref{fig:ridders-coverage} plots it against the actual error over the book's own
study: \num{\valRiddersFunctions} smooth functions with closed-form derivatives, \num{\valRiddersPointsPerFunction} random
points each in $[0.2, 3]$, the default initial step. A point above the diagonal is a case where the estimate
covered the error. Measured: \SI{\valRiddersCoverageFirst}{\percent} of the first derivatives and
\SI{\valRiddersCoverageSecond}{\percent} of the second are covered; the worst uncovered case has an error of
\num{\valRiddersWorstRatioFirst} and \num{\valRiddersWorstRatioSecond} times its estimate; and the median ratio of
error to estimate is \num{\valRiddersMedianRatioFirst} and \num{\valRiddersMedianRatioSecond}, so the typical
estimate is an overestimate by a factor of ten or so. Two things the figure does not say must be said: the
estimate is not a bound, and the study's scope is smooth functions, no noise, and for $f'$ no singularity within
the initial step (for the higher orders the larger initial step does reach a singularity at small $x$, which the
method survives by restarting, or, where the samples stay finite, by taking more levels,
\cref{sb:floor-coverage}). Outside that scope is the next section.

\section{When not to trust it}\label{sec:ridders-trust}

\paragraph{A non-finite value: the edge of a domain.} Near the edge of a domain Ridders' method still works, at a
higher cost; exactly at an edge it does not, and its one-sided forms (\cref{sec:ridders-onesided}) or the
one-sided stencils of \cref{sec:stencils-one-sided} are the tool. The initial step is $\scale / 32$ for $f'$ with $\scale = \max(\abs{x}, 1)$, which near an edge may reach
across it. The method then sees a NaN or an infinity, discards the tableau and restarts from half the failing
step, since every earlier level's interval contained the bad point too. Take $\log x$ at $x = 0.01$: the step
$1/32$ samples below zero, and so does $1/64$, so two restarts precede the tableau that converges:
\snippet{06-ridders-edge}{edge-restart}
Ridders returns \num{\valRiddersEdgeAtOneHundredthValue}, accurate to \num{\valRiddersEdgeAtOneHundredthRelError}
relative and covered by its estimate, after \num{\valRiddersEdgeAtOneHundredthEvals} evaluations against
\num{\valRiddersEdgeEvalsInterior} for $\log x$ away from its edge: four of the extra evaluations are the two
failed levels, and the rest a tableau that starts at $1/128$, nearly the size of $x$, and needs more levels to
converge. \Cref{fig:ridders-edge} measures how far this goes: each further halving of $x$ adds one failed level,
two evaluations; the level budget of \num{\valRiddersLevelCap} is used up from $x = 2^{-\num{\valRiddersEdgeCapExponent}}$;
ten correct digits hold down to $x = \num{\valRiddersEdgeTenDigits}$; a value is still returned down to
$x = \num{\valRiddersEdgeLimit}$, with the estimate risen to \num{\valRiddersEdgeLimitRelEstimate} of the value; and below that the result is
\code{non_finite}. As a rule, in its two readings: a value is returned as long as the nearest bad point of $f$ is
about $2^{-34} \scale$ away for $f'$, $2^{-31} \scale$ for $f''$ and $2^{-30} \scale$ for $f'''$ and $f''''$, whose
initial step is $\scale / 4$ and whose rules reach $2h$ (derived from the initial step and the level cap;
measured in \cref{sec:guarantees-limits}); for
ten digits keep it at least $2^{-\num{\valRiddersEdgeTenDigitsExponent}} \scale$ away (measured for $f'$). The stencil in the same example does not
fail, since its default step of $\epsmach^{1/3}$ never reaches zero, but its relative error is
\num{\valEdgeCentralAtOneHundredthRelError}: the truncation $h^2 / 3x^2$ with a step scaled by $\scale = 1$ rather
than by $x$, the small-$\abs{x}$ case of \cref{sec:rounding-scale}, which Ridders handles by adapting the step.

\begin{marginfigure}
\centering
% F15 fig:ridders-edge (Chapter 6, margin). Ridders for log x at x = 2^-k against k with the default initial step
% s/32 = 1/32 (DS10, data/ridders-edge.csv): above, the evaluations of fn (each halving of x costs one more restart,
% about two evaluations); below, the relative error, which stays near rounding until the level budget runs short.
% The three measured limits are marked from the generated values: the budget of 64 evaluations is first exhausted
% at k = \valRiddersEdgeCapExponent, ten correct digits last hold at k = \valRiddersEdgeTenDigitsExponent, and a
% value is last returned at k = \valRiddersEdgeLimitExponent; beyond it the result is non_finite (the nan rows break
% the curve). The last CSV row (x = 0.01, a non-integer k, E10) is dropped by the x expression.
\begin{tikzpicture}
\begin{groupplot}[group style={group size=1 by 2, vertical sep=3mm, xticklabels at=edge bottom},
  eqtufte, width=36mm, height=22mm, xmin=2, xmax=40, xtick={10,20,30,40}, unbounded coords=jump]
  \nextgroupplot[ymin=0, ymax=70, ytick={16,32,48,64}, ylabel={evaluations}]
  \addplot[black!75] table[col sep=comma, comment chars={\#},
    x expr={\thisrow{k} == floor(\thisrow{k}) ? \thisrow{k} : nan}, y=evaluations] {data/ridders-edge.csv};
  \draw[muted, thin] (axis cs:\valRiddersEdgeCapExponent, 0) -- (axis cs:\valRiddersEdgeCapExponent, 70);
  \draw[muted, thin] (axis cs:\valRiddersEdgeTenDigitsExponent, 0) -- (axis cs:\valRiddersEdgeTenDigitsExponent, 70);
  \draw[accent, thin] (axis cs:\valRiddersEdgeLimitExponent, 0) -- (axis cs:\valRiddersEdgeLimitExponent, 70);
  \node[muted, font=\tiny, anchor=south east, inner sep=1pt] at (axis cs:\valRiddersEdgeCapExponent, 70) {cap};
  \node[muted, font=\tiny, anchor=south, inner sep=1pt, yshift=1.1em] at (axis cs:\valRiddersEdgeTenDigitsExponent, 70) {10 digits};
  \node[accent, font=\tiny, anchor=south west, inner sep=1pt] at (axis cs:\valRiddersEdgeLimitExponent, 70) {last value};
  \nextgroupplot[ymode=log, log basis y=10, ymin=1e-16, ymax=1, ytick={1e-15,1e-10,1e-5,1},
    ylabel={relative error}, xlabel={$k$ ($x = 2^{-k}$)}]
  \addplot[black!75] table[col sep=comma, comment chars={\#},
    x expr={\thisrow{k} == floor(\thisrow{k}) ? \thisrow{k} : nan}, y=rel_error] {data/ridders-edge.csv};
  \draw[muted, thin] (axis cs:\valRiddersEdgeCapExponent, 1e-16) -- (axis cs:\valRiddersEdgeCapExponent, 1);
  \draw[muted, thin] (axis cs:\valRiddersEdgeTenDigitsExponent, 1e-16) -- (axis cs:\valRiddersEdgeTenDigitsExponent, 1);
  \draw[accent, thin] (axis cs:\valRiddersEdgeLimitExponent, 1e-16) -- (axis cs:\valRiddersEdgeLimitExponent, 1);
  \node[accent, font=\tiny, anchor=north west, inner sep=1pt] at (axis cs:\valRiddersEdgeLimitExponent, 1e-12) {non-finite};
\end{groupplot}
\end{tikzpicture}

\caption{Evaluations and relative error of Ridders for $\log x$ at $x = 2^{-k}$ against $k$: each halving of $x$
costs about two evaluations until the budget is used up at $k = \num{\valRiddersEdgeCapExponent}$; ten correct
digits hold to $k = \num{\valRiddersEdgeTenDigitsExponent}$, digits are lost from there, a value is returned to
$k = \num{\valRiddersEdgeLimitExponent}$ and the result is non-finite beyond.\label{fig:ridders-edge}}
\end{marginfigure}

\paragraph{A pole between the samples, or aliasing.} The samples can be finite and still lie: a pole of $f$
between two sample points gives finite differences that converge to nothing, and a function that oscillates faster
than the initial step can vanish at every sample. The method has a test for a result that has at most two correct
bits by its own estimate while that estimate is far above rounding, and restarts at half the smallest step it tried
when it fires (\cref{sb:aliasing}). That restart descends as the one after a non-finite value does, so a pole
close to $x$ is resolved: at the tests' poles, $10^{-3}$ ahead of 0 and $10^{-4}$ behind it, far inside the
limits above, \num{\valPoleCovered} of the \num{\valPoleCases} results for $f'$ to $f''''$ are covered, from
\num{\valPoleRelErrorMin} to \num{\valPoleRelErrorMax} off. Much closer, a pole gives no \code{non_finite} but a
wrong value, flagged by its estimate in every case measured (\cref{sec:guarantees-limits}): test \code{abs_error}.

\paragraph{A function less accurate than $\epsmach$.} The floor assumes that \code{fn} is accurate to about
$\epsmach$. A function that carries the noise $\noise$ of an iterative solver (\cref{sec:rounding-noise}) breaks
that assumption, and the estimate becomes optimistic: measured over the study's \num{\valNoiseRiddersCases} random
points of $\mathrm{e}^x$ with the values perturbed by $\noise = 10^{-8}$, the coverage falls to
\SI{\valNoiseEightRiddersCoverage}{\percent} with a worst ratio of \num{\valNoiseEightRiddersWorstRatio}
(\cref{fig:noise}). Pass an initial step on the scale of the noise, and do not read the estimate as a bound.

\paragraph{A jump or a kink inside the span.} A range break within $h_0$ of $x$ does not produce a non-finite
value, and Ridders returns a value for the wrong function. Where $f$ itself jumps, it also returns a large
estimate, since the base differences then converge like $\order(h)$ rather than $\order(h^2)$ and the entries keep
changing: on the nitrogen example's own range break the stencils return errors of order one in silence, while Ridders'
relative estimate is \num{\valKinkInsideRiddersRelEstimate}, ten orders of magnitude above the
\num{\valKinkOutsideRiddersRelEstimate} it reports \SI{\valKinkOutsideDelta}{\kelvin} from the break
(\cref{sec:using-pitfalls}, \cref{fig:kink}). The estimate flags a jump; the stencil does not. A kink, a jump in
$f'$ with $f$ continuous, at $x$ itself shows no symptom at all: the central difference of $c \abs{t}$ is exactly
zero at every step, so every central rule returns the mean of the two one-sided slopes plus the smooth part's own error,
Ridders' tableau converges to that mean with its estimate at the floor, and nothing in the samples can show it;
know the branch boundaries and differentiate the branch you mean (\cref{sec:using-pitfalls}).

\section{Cost}\label{sec:ridders-cost}

\begin{marginfigure}
\centering
% F14b fig:ridders-cost (Chapter 6, margin). Histogram of the evaluations of fn per call over the random study
% (DS9, data/ridders-coverage-N1.csv and -N2.csv): N = 1 (dark, even counts: 2 per level) and N = 2 (muted, odd
% counts: 2 per level plus fn(x)). Unit bins centred on the integers, so the two series interleave. The tail of
% N = 2 at 19 and 21 is the slower convergence of several functions; the 11 restarts of log and sqrt below x = 0.25 are a small part of the last bar (R1-decision T7).
\begin{tikzpicture}
\begin{axis}[eqtufte, width=38mm, height=30mm,
  xmin=5.5, xmax=22.5, ymin=0, ymax=2100,
  xtick={6,10,14,18,22}, ytick={0,1000,2000}, scaled y ticks=false,
  xlabel={evaluations of \code{fn}}, ylabel={cases}]
  \addplot[hist={bins=17, data min=5.5, data max=22.5}, fill=black!30, draw=none]
    table[col sep=comma, comment chars={\#}, y=evaluations] {data/ridders-coverage-N2.csv};
  \addplot[hist={bins=17, data min=5.5, data max=22.5}, fill=black!75, draw=none]
    table[col sep=comma, comment chars={\#}, y=evaluations] {data/ridders-coverage-N1.csv};
  \node[font=\tiny, anchor=south, inner sep=1pt] at (axis cs:14.5, 1900) {$\dorder = 1$};
  \node[font=\tiny, anchor=south, inner sep=1pt, text=black!55] at (axis cs:17.5, 1580) {$\dorder = 2$};
\end{axis}
\end{tikzpicture}

\caption{Evaluations of \code{fn} per call over the random study: $\dorder = 1$ dark, $\dorder = 2$ muted. The
tail of $\dorder = 2$ is the slower convergence of several functions, not the restarts, which are a small part of
its last bar (\cref{sec:ridders-cost}).\label{fig:ridders-cost}}
\end{marginfigure}

A level costs 2 evaluations of \code{fn} for $f'$ and $f''$ and 4 for $f'''$ and $f''''$, plus $f(x)$ once per
call for the even orders. Measured over the random study (\cref{fig:ridders-cost}), a first derivative takes a
median of \num{\valRiddersEvalsMedianFirst} evaluations, between \num{\valRiddersEvalsMinFirst} and
\num{\valRiddersEvalsMaxFirst}, and a second \num{\valRiddersEvalsMedianSecond}, between \num{\valRiddersEvalsMinSecond} and
\num{\valRiddersEvalsMaxSecond}. The tail of the second-derivative histogram, the \num{\valRiddersTailCasesSecond}
cases of \num{\valRiddersTailMinEvalsSecond} evaluations or more, is the slower convergence of functions whose
higher derivatives are large relative to $f''$: \valRiddersTailFunctionsSecond; the \num{\valRiddersTailRestartsSecond}
restarts of $\log x$ and $\sqrt{x}$ below $x = 0.25$, where the initial step $\scale / 4$ reaches below zero, all
take the most evaluations and are a small part of the last bar. The third and fourth derivatives take about twice
as many (\cref{tab:cost}). The caps are \num{\valRiddersLevelCap} levels in all, restarts included, which is at
most \num{\valRiddersMaxEvalsFirst}, \num{\valRiddersMaxEvalsSecond}, \num{\valRiddersMaxEvalsThird} and
\num{\valRiddersMaxEvalsFourth} evaluations for the four orders; \cref{sec:guarantees-ridders} collects the
statistics for all four. Against a stencil's two to six evaluations, Ridders takes 2.3 to 7 times as many for
$f'$, and for them returns about three more digits than \code{Central<2>} (medians \num{\valExpCentralTwoRelError}
and \num{\valExpRiddersFirstRelError} on $\mathrm{e}^x$: \num{\valHeadlineFirstDigits} against
\num{\valHeadlineRiddersFirstDigits} digits), as many as \code{Central<4>} (\num{\valExpCentralFourRelError}), none
over \code{Central<6>} (\num{\valExpCentralSixRelError}, \cref{tab:accuracy-matrix}), and an estimate.

\section{Stencil or Ridders}\label{sec:ridders-choice}

Why does only Ridders estimate its error? A stencil could bound its truncation error only with extra evaluations,
at a second step or a higher order, and its rounding part alone is smallest exactly where truncation is largest;
an estimate that costs nothing would be worth nothing. So the library leaves a stencil's \code{abs_error} empty,
and a code that needs an estimate asks for Ridders. The choice follows from that. Inside an iteration, where the
derivative is recomputed many times and a Newton step tolerates an approximate slope, a stencil with the default
step is the right tool: two or four evaluations, a relative error of about \num{\valExpCentralTwoRelError} or
\num{\valExpCentralFourRelError}, no decisions. For a value that is reported, compared or published, Ridders:
several times the evaluations, about three more digits than \code{Central<2>} (\num{\valHeadlineRiddersFirstDigits}
against \num{\valHeadlineFirstDigits}) and as many as \code{Central<4>}, and
a number that says how far to trust it (a stencil's result can be checked as in \cref{sec:using-check}).

\section{One-sided Ridders}\label{sec:ridders-onesided}

\begin{marginfigure}
\centering
% F19 fig:ridders-onesided-reach (Chapter 6, margin; Phase 5). The reach of ForwardRidders from its default initial
% step s/32 = 1/32: log x at x = 2^-k against k (DS23, data/forward-ridders-reach.csv). Relative error (solid) and
% the estimate over the value returned (dashed, abs_error / |value|, the quantity the flag compares with 1/4; the
% thin line at 1/4, R3-numerics-04, R3-evidence-03) for f' (ink) and f'''' (accent), and central Ridders' relative
% error for f' (muted,
% dotted; its samples cross 0, where log is NaN, and its restarts follow that NaN; no value beyond
% k = \valRiddersEdgeLimitExponent, the nan rows break the curve). Beyond its reach the one-sided result tends to 0,
% so the error tends to 1 and the estimate over the value rises above the line at 1/4: the result is flagged. The
% thin muted vertical line is the first k
% at which f' spends all its evaluations (\valForwardRiddersReachCapExponentFirst); the open circle is the one
% distance where f' is neither covered nor flagged (k = \valForwardRiddersReachSilentExponentFirst, D-18). The x axis
% carries k and, below it, the order of magnitude of x = 2^-k, as in F15.
\begin{tikzpicture}
\begin{axis}[eqtufte, width=38mm, height=56mm, ymode=log, log basis y=10,
  xmin=2, xmax=45, xtick={10,20,30,40},
  xticklabels={10\\$\approx 10^{-3}$, 20\\$\approx 10^{-6}$, 30\\$\approx 10^{-9}$, 40\\$\approx 10^{-12}$},
  xticklabel style={align=center, font=\sffamily\tiny},
  ymin=1e-16, ymax=10, ytick={1e-15,1e-10,1e-5,1},
  xlabel={$k$, and $x = 2^{-k}$}, ylabel={relative error, estimate}, unbounded coords=jump]
  \draw[muted, thin] (axis cs:\valForwardRiddersReachCapExponentFirst, 1e-16)
    -- (axis cs:\valForwardRiddersReachCapExponentFirst, 10);
  % central Ridders, f'
  \addplot[muted, densely dotted, thick]
    table[col sep=comma, comment chars={\#}, x=k, y=rid1_error] {data/forward-ridders-reach.csv};
  % ForwardRidders, f' and f''''
  \addplot[black!80]
    table[col sep=comma, comment chars={\#}, x=k, y=fr1_error] {data/forward-ridders-reach.csv};
  \addplot[black!80, dashed]
    table[col sep=comma, comment chars={\#}, x=k, y=fr1_estimate_value] {data/forward-ridders-reach.csv};
  \addplot[accent]
    table[col sep=comma, comment chars={\#}, x=k, y=fr4_error] {data/forward-ridders-reach.csv};
  \addplot[accent, dashed]
    table[col sep=comma, comment chars={\#}, x=k, y=fr4_estimate_value] {data/forward-ridders-reach.csv};
  % the flag: an estimate of at least a quarter of the value returned
  \draw[muted, thin] (axis cs:2, 0.25) -- (axis cs:45, 0.25);
  \node[muted, font=\tiny, anchor=south west, inner sep=1pt] at (axis cs:3, 0.25) {$1/4$};
  % the one f' result neither covered nor flagged
  \draw[black!80, thin] (axis cs:\valForwardRiddersReachSilentExponentFirst, \valForwardRiddersReachSilentRelErrorFirst)
    circle[radius=1.6pt];
  \node[font=\tiny, anchor=north west, inner sep=1pt] at (axis cs:3, 1e-12) {$f'$};
  \node[accent, font=\tiny, anchor=south west, inner sep=1pt] at (axis cs:3, 1e-4) {$f''''$};
  \node[muted, font=\tiny, anchor=north west, inner sep=1pt, align=left] at (axis cs:3, 1e-15) {central $f'$};
  \node[muted, font=\tiny, anchor=north east, inner sep=1pt, align=right]
    at (axis cs:\valForwardRiddersReachCapExponentFirst, 5) {levels\\used up};
\end{axis}
\end{tikzpicture}

\caption{The one-sided form for $\log x$ at $x = 2^{-k}$ from its default initial step: relative error (solid),
and estimate over the value returned (dashed), for $f'$ (dark) and $f''''$ (accent); central Ridders' error
for $f'$ (dotted). Each keeps
its accuracy for some halvings after all its evaluations are spent (the line, for $f'$); then the one-sided error
rises towards 1, covered for a while; further out, to the end of the study, the dashed curve is above the line
at $1/4$: the result is flagged. At one distance, for $f'$ (the circle), it is neither covered nor flagged. Central Ridders is more accurate and returns no value beyond
$k = \num{\valRiddersEdgeLimitExponent}$.\label{fig:ridders-onesided-reach}}
\end{marginfigure}

Where $f$ is undefined or different on one side of $x$, at a domain edge, at the end of a correlation's range, or
for the right-hand or left-hand slope at a kink, Ridders' method has one-sided forms. \code{ForwardRidders}
extrapolates the \code{Forward<1>} rule of order $\dorder$, for $f'$ the quotient $(f(x + h_i) - f(x)) / h_i$
(over the actual distance and offsets, \cref{sb:actual-distance,sb:actual-offsets}), and
evaluates \code{fn} at $x$ and above only, restarts included; \code{BackwardRidders} applies it to $f(-t)$ at
$-x$ and negates the value for odd $\dorder$, so it samples at $x$ and below, with the same accuracy, cost and
limits.

The forward rule's error has every power of the step, $b_1 h + b_2 h^2 + \cdots$, so column $j$ removes the term
in $h^j$ and divides by $2^j - 1$ instead of $4^j - 1$: one order per level instead of two, the general case
of Burden et al.\cite{burden2016} (Section 4.2; \cref{sec:derivations-richardson}). The smaller divisors amplify rounding more, so the floor of an entry is carried
through the tableau, $r_{i,j} = (2^j r_{i,j-1} + r_{i-1,j-1}) / (2^j - 1)$ from $r_{i,0} = r_i$, rather than taken
from its level; the rest, restarts and cap included, is Ridders' method as described. The default initial step is
$\scale / \num{\valForwardRiddersDivisor}$ for every order, chosen on a prototype sweep (the design record). A call evaluates
$f(x)$ once and $\dorder$ points per level, at most \num{\valForwardRiddersMaxEvalsFirst} to
\num{\valForwardRiddersMaxEvalsFourth} times, typically \num{\valForwardRiddersEvalsMedianFirst} for $f'$ and
\num{\valForwardRiddersEvalsMedianFourth} for $f''''$.

The side has a price. On the study of \cref{fig:ridders-coverage} (the same smooth functions and points, finite
above every $x$, no noise, double) the median relative error runs from \num{\valForwardRiddersMedianRelErrorFirst}
for $f'$ to \num{\valForwardRiddersMedianRelErrorFourth} for $f''''$, \num{\valForwardRiddersAccuracyFactorFirst}
to \num{\valForwardRiddersAccuracyFactorFourth} times Ridders' on the same cases; the estimate covers
\SI{\valForwardRiddersCoverageMin}{\percent} to \SI{\valForwardRiddersCoverageMax}{\percent} of them, the
worst error \num{\valForwardRiddersWorstRatioMax} times its estimate. A pole just ahead of $x$, inside the span
the rule samples, can defeat it for $f''''$, unflagged, well inside the reach below (\cref{sec:guarantees-ridders}). Where $f$ is smooth on both sides, use Ridders,
also near an edge, where the non-finite values guide its restarts. At the nitrogen fit's lower end, with
the wrapper of \cref{sec:stencils-one-sided}:
\needspace{6\baselineskip}%
\snippet{06-ridders-onesided}{onesided-edge}
Ridders returns \valEdgeRiddersStatus{} after \num{\valEdgeRiddersEvals} evaluations. The one-sided form, from
an initial step of about \SI{\valEdgeForwardRiddersInitialStep}{\kelvin}, returns $\cp$ \num{\valEdgeForwardRiddersRelError} off,
estimate \num{\valEdgeForwardRiddersRelEstimate}, after \num{\valEdgeForwardRiddersEvals} evaluations, and
$\dd \cp / \dd T$ \num{\valEdgeForwardRiddersSecondRelError} off, \num{\valEdgeForwardRiddersSecondRelEstimate},
after \num{\valEdgeForwardRiddersSecondEvals}, all relative.

Its reach is limited the other way: where $f$ varies on a scale far below the initial step, as $\log x$ near 0
does, a one-sided base sees no non-finite value, only untrusted results, and restarts at half the smallest step
it tried; beyond its reach the levels run out first (\cref{fig:ridders-onesided-reach}). On $\log x$ at $x = 2^{-k}$, $f'$ has ten
digits to $k = \num{\valForwardRiddersReachTenDigitsExponentFirst}$; at the documented reach, $k = 30$, $f'$ is
\num{\valForwardRiddersReachRelErrorAtThirtyFirst} off, $f''''$ \num{\valForwardRiddersReachRelErrorAtThirtyFourth}. The estimate covers the growing error to
$k = \num{\valForwardRiddersReachCoveredExponentFirst}$ for $f'$ (\num{\valForwardRiddersReachCoveredExponentFourth}
for $f''''$), and from $k = \num{\valForwardRiddersReachFlaggedExponentFirst}$
(\num{\valForwardRiddersReachFlaggedExponentFourth}) to the end of the study, $k = \num{\valForwardRiddersReachLastK}$,
the wrong result is flagged by an estimate of at least a quarter of its value. At
$k = \num{\valForwardRiddersReachSilentExponentFirst}$ between them, $f'$ is
\num{\valForwardRiddersReachSilentRelErrorFirst} off with an estimate of
\num{\valForwardRiddersReachSilentRelEstimateFirst}: neither covered nor flagged. As a rule, keep $f$'s
nearest feature at least about $2^{-30} \scale$ from $x$, or pass an initial step on its scale.

\begin{sidebar}[label=sb:difference-form]{Burden's difference form of the tableau}
\Cref{eq:ridders-richardson} generalises to the weighted combination
$(4^j \Aij{i}{j-1} - \Aij{i-1}{j-1}) / (4^j - 1)$ for column $j$. Burden et al.\cite{burden2016} (Section 4.2) write
the equal correction form,
\[
  \Aij{i}{j} = \Aij{i}{j-1} + \frac{\Aij{i}{j-1} - \Aij{i-1}{j-1}}{4^j - 1},
\]
and the library computes that form, for three reasons. When the two parents agree exactly, the correction is 0
and the entry equals them bit for bit, which is what the stop on exact agreement relies on; the product form
could round. The correction is a small quantity divided by a number above 1, so it cannot overflow where the
entries themselves are finite, whereas $4^j \Aij{i}{j-1}$ can exceed the largest finite value in float from
$j = 7$ on for a derivative near $2^{115}$; the library still checks the one subtraction and the one addition
before forming them, since two finite entries of opposite sign can differ by more than the largest value, and
treats a failure as a failed level. And only the previous row is needed: the tableau is two arrays of
\num{\valRiddersLevelCap} entries, filled from \code{T(0)} so that the real type need not be default-constructible,
and swapped after each level. The factors $4^j$ are built by repeated multiplication in the real type and are
exact in binary; $4^j - 1$ needs $2j$ bits and rounds once $2j$ exceeds the significand, from $j = 13$ in float and
$j = 27$ in double (far beyond the cap in a 50-digit type), by at most half an epsilon, which does not matter.
\end{sidebar}

\begin{sidebar}[label=sb:floor-coverage]{How Ridders stops, and how often its estimate holds}
\emph{The rules}, as the code has them. Level $i$ is the \code{Central<2>} rule for the order at $h_i = h_0 /
2^i$; before it is evaluated the method checks that $h_i$ is positive, that $x \pm h_i$ differ from $x$ (and for
$f'''$ and $f''''$ that $x \pm 2 h_i$ differ from $x \pm h_i$), and that no point can overflow; at the resolution
of $x$ the tableau simply ends. Each extrapolated entry gets the estimate $e_{i,j}$ of \cref{eq:ridders-estimate}
with the floor $r_i$ of its level; for $f'$ the floor is the rounding of the two samples over their distance, for
$f''$ the same from the three samples and the two offsets, and for $f'''$ and $f''''$ the rounding of the samples
plus the rounding of the points, $\tfrac12 \epsmach (\abs{x} + 2h)\, \abs{f'} \sum_k \abs{c_k} / (\den h^{\dorder})$
with the integer numerators $c_k$ over the denominator $\den$ and $f'$ taken from the samples at $x \pm h$, since
those rules divide by the nominal step. The best entry is the one with the smallest estimate since the last
restart. The tableau stops when an entry equals its parent exactly (nothing is left to learn; that entry is best,
with the floor as its estimate), when two consecutive levels bring no smaller estimate, at the resolution of $x$,
or when the level budget runs out. The result is the best entry with $\errest = \max(e_{\mathrm{best}},
\abs{A_{\mathrm{best}} - \Aij{i}{i}})$ for the latest diagonal entry; it needs at least two successful levels
since the last restart.

\emph{Restarts.} A level fails if \code{fn} is not finite at a point, if a difference, quotient or tableau entry
would overflow, or if the points would; the tableau, its best entry and the stall counter are discarded and a
new tableau starts from half the failing step. A result is untrusted if $\errest \ge \abs{A_{\mathrm{best}}} /
4$, at most two correct bits, and $\errest / 64 > r_{\mathrm{best}}$, more than rounding explains; the method then
restarts at half the smallest step it tried, as after a failed level, and keeps the untrusted result as
a fallback, returned if the budget runs out before a trusted one appears. A tableau that ended at the resolution of $x$ is returned as it is, since no restart
could go lower; a derivative that is exactly 0 has an estimate near its floor and is not affected. The budget is
\num{\valRiddersLevelCap} levels in all; for the even orders $f(x)$ is evaluated once and kept across restarts, and if
it is not finite the result is \code{non_finite} after that one evaluation.

\emph{Where the constants come from.} The ratio 2, the two-level stall, the factor 1 on the floor, the factors 4
and 64 of the untrusted test and the cap of 32 were chosen with a prototype on random cases, and the design record
in the repository, \code{docs/deriv/PHASE3_API.md} to \code{docs/deriv/PHASE5_API.md}, gives the evidence for each: without
the floor the estimate fell below the actual error in more than half of the prototype's cases; a one-level stall
let a Runge-function case stop early; an untrusted test at $\abs{A}$ (factor 1) instead of $\abs{A} / 4$ left
poles between samples uncovered; factors between 16 and 256 over the floor gave identical results. The book's own
evidence is the study of \cref{fig:ridders-coverage}: \num{\valRiddersFunctions} functions ($\mathrm{e}^x$,
$\sin x$, $\log x$, $\sqrt{x}$, $1/x$, the Runge function $1/(1 + 25x^2)$, $x^5$, $\mathrm{e}^{-x^2}$ and
$\tanh x$) at \num{\valRiddersPointsPerFunction} random points each in $[0.2, 3]$, with the default initial step,
no noise and no singularity within $h_0$ for $f'$. Measured: coverage \SI{\valRiddersCoverageFirst}{\percent}
for $f'$ (\num{\valRiddersUncoveredFirst} of \num{\valRiddersCases} cases uncovered) and
\SI{\valRiddersCoverageSecond}{\percent} for $f''$ (\num{\valRiddersUncoveredSecond}); worst ratios of error to
estimate \num{\valRiddersWorstRatioFirst} and \num{\valRiddersWorstRatioSecond}; median ratios
\num{\valRiddersMedianRatioFirst} and \num{\valRiddersMedianRatioSecond}. The uncovered cases have a shape in
common: for $f'$, \num{\valRiddersUncoveredMainCountFirst} of the \num{\valRiddersUncoveredFirst} are
\valRiddersUncoveredMainFunctionFirst{} at $x$ between \num{\valRiddersUncoveredMainMinXFirst} and
\num{\valRiddersUncoveredMainMaxXFirst} (the other: \valRiddersUncoveredOthersFirst), and for $f''$ all
\num{\valRiddersUncoveredMainCountSecond} are \valRiddersUncoveredMainFunctionSecond{} at $x$ between
\num{\valRiddersUncoveredMainMinXSecond} and \num{\valRiddersUncoveredMainMaxXSecond}, where the derivatives grow
like $(2x)^k$ and the first levels are not yet in the asymptotic regime, so the change from the parent undershoots
the entry's error; a function whose derivatives grow geometrically on the scale of $h_0$ is where the estimate is
tightest, and a smaller initial step usually restores the margin: a quarter of the default initial step covers
\num{\valRiddersMarginQuarterCovered} of the study's \num{\valRiddersMarginCases} uncovered cases (all four
orders), a sixty-fourth \num{\valRiddersMarginSixtyFourthCovered}. There were no restarts for $f'$ and \num{\valRiddersRestartedSecond} for $f''$, the
cases of $\log x$ and $\sqrt{x}$ below $x = 0.25$, where $h_0 = \scale / 4$ reaches below zero and the first levels
are NaN. $1/x$ at the same points does not restart: its pole lies between finite samples, the first levels are
simply wrong, and the tableau takes more levels and converges from the later ones, accurately.
\end{sidebar}

\needspace{5\baselineskip}%
\begin{sidebar}[label=sb:aliasing]{Aliasing: when the samples lie}
Take $f(x) = \sin(64 \pi x)$ at $x = 0$, whose derivative is $64\pi$. The default initial step is $1/32$, so the
first level samples $\sin(\pm 2\pi)$ and the second $\sin(\pm \pi)$: every sample is zero up to rounding, the two
base differences nearly agree, and the first extrapolated entry is a confident 0 with a tiny estimate.
What saves the call is the sequence: the third level
samples $\sin(\pm \pi / 2) = \pm 1$, and the diagonal entries jump to numbers of order 100 while the best entry
stays near 0. The final check raises the estimate to the distance from the latest diagonal entry, the untrusted
test fires, and the method restarts at half the smallest step it tried; the second tableau converges to $64\pi$
(measured: \valAliasingResult{}, after \num{\valAliasingEvals} evaluations).
The design record names the case that cannot be caught: samples that are exactly periodic \emph{and} lie on a
polynomial, such as $\cos(64 \pi x)$ at 0 for $f''$, whose first samples are all exactly 1: two entries agree
exactly and the method stops with the floor as its estimate. That fools every sampling method;
the remedy is a step on the scale of the period.
\end{sidebar}
